% =====================================================================
%  4 Versionsverwaltung  (ab Phase 2)
%  Aufgabenstellung 1.2: „Der Link zum GitHub-Repository ist an geeigneter
%  Stelle in der Projektdokumentation zu ergänzen.“
% =====================================================================
\section{Versionsverwaltung und Repository}\label{sec:versionsverwaltung}
Quellcode, Dokumente, Diagrammquellen und alle Dateien, die zur Ausführung nötig sind, liegen gemeinsam im folgenden Git-Repository auf GitHub:

\url{\GitHubURL}

Das Repository ist öffentlich. Tutor:innen können es daher ohne Einladung und ohne eigenes Konto lesen, klonen oder als ZIP-Datei herunterladen. Die Aufgabenstellung der IU ist urheberrechtlich geschützt und liegt deshalb nicht im Repository. Zusätzlich enthält Anhang~\ref{anh:quellcode} den Quellcode der Anwendung in kopierfähiger Form.

Das Repository ist in vier Hauptordner gegliedert:
\begin{itemize}
  \item \texttt{app/}: die Anwendung mit Quellcode (\texttt{src/balanceos}), Tests (\texttt{tests}), Projektkonfiguration (\texttt{pyproject.toml}) und der Konfiguration für PyInstaller \parencite{pyinstaller}, das daraus die \texttt{.exe} baut. Die ausführbare Datei liegt in \texttt{app/release/BalanceOS.exe}. Die Datei \texttt{app/README.md} beschreibt Installation, Start und Tests.
  \item \texttt{docs/}: die \LaTeX-Quellen aller Dokumente, die TikZ-Quellen der Diagramme mit ihren PDF-Exporten und die gebauten Dokumente in \texttt{docs/pdf}. Der Ordner \texttt{docs/superpowers} enthält die Design- und Planungsnotizen, die bei der Arbeit mit dem \ac{KI}-Werkzeug je Phase entstanden sind.
  \item \texttt{abgabe/}: die eingereichten Abgabe-PDFs und Kurzzusammenfassungen je Phase.
  \item \texttt{tools/}: das Skript \texttt{build\_docs.py}, das Diagramme und Dokumente baut und die Abgabe-PDF einer Phase zusammenstellt.
\end{itemize}

Jede Phase wird auf einem eigenen Zweig bearbeitet, etwa in Phase~3 auf \texttt{phase3}. Am Ende der Phase wird dieser Zweig in \texttt{main} zusammengeführt. Nach der Einreichung markiert ein Tag den eingereichten Stand (\texttt{phase1-abgabe}, \texttt{phase2-abgabe}, \texttt{phase3-abgabe}). So lässt sich jede Abgabe später genau wiederherstellen. Commit-Nachrichten sind auf Deutsch verfasst. In der Implementierung beginnen sie mit dem Inkrement, etwa \enquote{I3: Regelwerk R01--R10 mit Begründung und Evidenzkennzeichnung}. Jedes Inkrement endet mit erfolgreich durchlaufenen Tests und wird auf GitHub hochgeladen. Commits, an denen das \ac{KI}-Werkzeug Claude Code beteiligt war, tragen eine \texttt{Co-Authored-By}-Zeile (siehe Abschnitt~\ref{sec:vorgehen}).
